\newpage

\appendix
\section{Extended Definitions and Proofs}

\noindent
\textbf{Full definition of the rules of Def.~\ref{def:side_effect} at page~\pageref{def:side_effect}.}
The complete definition of the rules for the side-effect of expressions and
commands is given in Fig.~\ref{fig:side_effect_full}.
%
\begin{figure}[b]
  \begin{align*}
    \sideeffect{c}&=\sideeffect{v}=\varnothing\\
    \sideeffect{\field{\exp}{f}}&=\sideeffect{\exp}\\
    \sideeffect{\aop{\exp_1}{\exp_2}}&=\sideeffect{\exp_1}\cup\sideeffect{\exp_2}\\
    \sideeffect{\new{C}{\exp_1,\ldots,\exp_n}}&=I(C.C)\cup\bigcup_{i=1}^n\sideeffect{\exp_i}\\
    \sideeffect{\size{\exp}}=\sideeffect{\last{\exp}}&=\sideeffect{\random{\exp}}=\sideeffect{\exp}
  \end{align*}
  \begin{align*}
    \sideeffect{\assign{v}{\exp}}&=\sideeffect{\exp}\\
    \sideeffect{\assign{\field{\exp_1}{f}}{\exp_2}}&=\sideeffect{\exp_2}\cup\sideeffect{\exp_2}\\
    \sideeffect{\require{\exp_1}{\exp_2}}&=\sideeffect{\receive{\exp_1}{\exp_2}}=\sideeffect{\exp_1}\cup\sideeffect{\exp_2}\\
    \sideeffect{\push{\exp_1}{\exp_2}}&=\sideeffect{\exp_1}\cup\sideeffect{\exp_2}\cup
    \{t\mid\exp_1\text{ has type }\<List<>t\<\text{>}>\}\\
    \sideeffect{\call{\exp_0}{m}{\exp_1,\ldots,\exp_n}}&=\bigcup_{j=1}^k I(C_j.m)\cup\bigcup_{i=0}^n\sideeffect{\exp_i}\\
    \noalign{\centering\text{where $C_1.m\cdots C_k.m$ are the methods possibly called here}}
    \sideeffect{\nop}&=\varnothing\\
    \sideeffect{\ifetwo{\exp_1}{\exp_2}{\com_1}{\com_2}}&=\sideeffect{\exp_1}\cup\sideeffect{\exp_2}\cup\sideeffect{\com_1}\cup\sideeffect{\com_2}\\
    \sideeffect{\for{v}{\exp}{\com}}&=\sideeffect{\exp}\cup\sideeffect{\com}\\
    \sideeffect{\sequence{\com_1}{\com_2}}&=\sideeffect{\com_1}\cup\sideeffect{\com_2}.
  \end{align*}
  \caption{The side-effect rules. They use an interpretation $I$ for the methods.}
  \label{fig:side_effect_full}
\end{figure}

\mbox{}\\
\noindent
\textbf{Proof of Prop.~\ref{prop:termination} at page~\pageref{prop:termination}.}
  By~\cite{CousotC77}, one must prove that $\nabla$ is an upper bound operator and that
  its application to any ascending chain stabilizes. The upper bound property holds since
  for every $(\size{\pi}\le n)\in\nabla(\|bo|_1,\|bo|_2)$ it is
  $(\size{\pi}\le n)\in\|bo|_1$ and $(\size{\pi}\le n)\in\|bo|_2$.
  Therefore, $\|bo|_1\sqsubseteq\nabla(\|bo|_1,\|bo|_2)$ and
  $\|bo|_2\sqsubseteq\nabla(\|bo|_1,\|bo|_2)$.
  Consider an ascending chain $\|bo|_0\sqsubseteq\|bo|_1\sqsubseteq\cdots\sqsubseteq
  \|bo|_k\sqsubseteq\cdots$ now. The application of the widening yields the sequence
  $\|bo|'_0=\|bo|_0$, $\|bo|'_{k+1}=(\|bo|'_k\sqcup\|bo|_{k+1})\nabla\|bo|'_k$ for $k\ge 0$.
  By Def.~\ref{def:widening}, it is
  $\|bo|'_0\supseteq\cdots\supseteq\|bo|'_{k}\supseteq\|bo|'_{k+1}\supseteq\cdots$.
  If, by contradiction, that sequence were not converging, then one can extract an
  infinite strictly decreasing subsequence
  $\|bo|''_0\supset\cdots\supset\|bo|''_{k}\supset\|bo|''_{k+1}\supset\cdots$, which is
  impossible since $\|bo|_0$ and $\|bo|_0'$ are assumed to be finite and
  consequently also $\|bo|_0''\subseteq\|bo|_0$ is finite.
  Note that the fact that $\nabla$ is a widening operator on \emph{finite} elements
  of $\bo$ is sufficient for termination, since the rules in
  Fig.~\ref{fig:bo_derivations_expressions}, \ref{fig:bo_derivations_commands1}
  and~\ref{fig:bo_derivations_commands2}
  lead to finite states if started from $\varnothing$, as one can be easily prove by induction.
  \qed

